\chapter{Twelve weeks to a decision}
\label{ch:implementation}

At the end of twelve weeks, the sponsor should be able to start with a real
authoring brief, watch one document move from plan to pre-human release, inspect
one measured feasibility case, and price the next decision. The build must test
the product proposition without leaving a stranded platform if the case fails.
Twelve weeks is a feasibility horizon, not a promise of market readiness.

Figure~\ref{fig:twelve-week-schedule} is the canonical calendar for this
proposal. Later work-package descriptions refer to these same intervals rather
than introducing a second schedule.

\hvFigureSchedule

\section{The twelve-week sequence}

\begin{table}[H]
\centering
\small
\begin{tabularx}{\textwidth}{@{}p{0.12\textwidth}p{0.25\textwidth}L p{0.24\textwidth}@{}}
\toprule
Weeks & Work & Deliverable & Decision consequence \\
\midrule
1--2 & Freeze the authoring brief, threat model, evidence boundary, named
exemplars, known vocabulary, rights manifest, corpus split, schemas, and
protected-object fixtures. & One versioned brief, contract, rights record,
reader rubric, and parser test set. & Stop if the decision, reader, evidence
boundary, or failure cost cannot be named. \\
3--6 & Implement and replay the protected core: the common document
representation, source locations, render checks, protected spans, citation
fixtures, canonical build, versioned records, and safe command contract. &
A second operator can reproduce the LaTeX source map and protected comparison.
& At week six, narrow to the protected review utility if replay fails. \\
7--9 & Conditional on the core gate, add the blueprint, bounded writer,
integrity, argument, writing, evidence, and reconstruction preflight, mutation
harness, and bounded repair. & A fail-closed preflight record and an extension
that passes calibration and held-out replay. & Keep model-dependent gates in
abstention if their false negatives or burden are not measured. \\
10--12 & Conditional on the same gate, integrate author dispositions, release
gate, reader packet, privacy controls, interaction log, and the live feasibility
case. & One document can move from approved brief to a reproducible
review-ready packet, followed by a timed reader task, without exposing private
workflow language. & Stop if a failed gate can emit a packet or the packet
cannot be reproduced. \\
\bottomrule
\end{tabularx}
\caption{The implementation schedule produces evidence that changes the next
decision, not a sequence of status reports.}
\label{tab:implementation-sequence}
\end{table}

The stages overlap where the dependency is clear. A trustworthy diff needs the
parser and protected representation first. The author packet needs the diff
and findings. The live reader needs the packet. A package benchmark can inform
those steps, but it cannot replace the end-to-end case.

\section{People and resource envelope}

The repository does not contain local loaded rates or document volume, so the
plan uses person-weeks. The sponsor can substitute those rates before making a
funding decision.

\begin{table}[H]
\centering
\small
\begin{tabularx}{\textwidth}{@{}p{0.24\textwidth}p{0.16\textwidth}L@{}}
\toprule
Role & Planning effort & Responsibility \\
\midrule
Product engineer & 12 person-weeks & Brief and blueprint schema, source
representation, bounded writer contract, protected comparison, diagnostics,
release gate, interfaces, build, and privacy controls. \\
Evaluation lead & 6 person-weeks & Corpus, mutation harness, critic calibration,
reader rubric, comparator, feasibility measurements, and priced comparative
design. \\
Domain editor and writing reviewer & 3 person-weeks & Protected-claim inventory,
technical-genre interpretation, blueprint review, and resolution of disputed
scientific meaning. \\
Security or policy owner & 1 person-week plus incident availability & Threat
fixtures, network policy, retention and disclosure boundary, and security veto.
\\
Independent coding reviewer & 1 person-week & Blind coding of reader answers
and finding labels; freezes decisions before seeing workflow assignment. \\
Independent reader(s) & 2 sessions plus preparation & Reads by people unfamiliar with the project and
recorded acceptance or requested changes. \\
Sponsor or project owner & 2 decision sessions & Select the live case, approve
the comparator, and choose proceed, revise, or stop. \\
\bottomrule
\end{tabularx}
\caption{A planning shape, not a dollar budget. The security or policy owner,
domain editor, and evaluation lead are explicit roles; one person may hold more
than one role in a small pilot, but the decisions remain distinct.}
\label{tab:resource-envelope}
\end{table}

The cost worksheet multiplies these person-weeks and reader sessions by local
loaded rates, then adds approved compute, data, and external-review costs. Those
observed or quoted inputs enter $C$ in Equation~\eqref{eq:value}. A short
calendar should not be mistaken for a cheap project, but another displayed
budget formula would add no information before the rates exist.

The protected-source core is the week-six checkpoint. Weeks 1--6 must deliver a
replayable source map, protected comparison, and safe command contract before the
bounded writer or model critics consume the remaining horizon. If that core does
not pass with a second operator, the sponsor narrows the build to the protected
review utility. The full authoring extension is conditional, not an implicit
promise that one engineer can productionize every command in twelve weeks.

\section{Deliverables at the end of the build}

The sponsor should receive six reader-visible results:

\begin{enumerate}
\item a local-first command entry point for LaTeX, from brief to release;
\item a source map, claim map, dependency graph, and protected-object report
      that can be inspected without trusting a model;
\item a bounded writer adapter plus a small, adjudicated corpus with tuning,
      mutation, and held-out partitions;
\item the HPO density calibration fixture, including positive ordering cases
      and adjudicated false alarms;
\item one end-to-end, fail-closed feasibility packet from a live document;
\item a measured account of first-draft release rate, repair cycles, author
      burden, reader time, abstentions, and meaning-preservation checks; and
\item a costed comparative-study design with a proceed, revise, or stop
      recommendation.
\end{enumerate}

The implementation may also produce machine-readable logs, package manifests,
and test fixtures. Those are supporting records. They should not become a
second narrative or be pasted into the manuscript as evidence of value.

\section{Dependencies and external decisions}

Four external decisions remain with the owner; they are listed on the next
page.

\begin{description}
\item[Live document.] The owner must approve a document whose failure cost is
  meaningful and whose reader can participate without compromising
  confidentiality.
\item[Privacy boundary.] The owner must decide whether any remote model call is
  permitted and what retention class applies to excerpts and rendered pages.
\item[Comparator.] The evaluation lead and owner must describe the incumbent
  workflow before the Humanvoice case begins.
\item[Evidence completion.] A librarian or domain-method reviewer must complete
  the outstanding literature and package checks before a strong scholarly or
  efficacy claim is circulated.
\end{description}

These are real dependencies, not forms to be completed after the build. If one
cannot be resolved, the project can still develop parser fixtures or a
synthetic review packet, but it must not present that work as a live efficacy
result.

\section{Risks with owners and stop triggers}

Human judgment is not immune to gaming, and a local system is not private by
mere intention. The operating plan therefore gives each important failure an
owner, an early signal, and a point at which the team stops or removes a
feature.

\begin{longtable}{@{}>{\raggedright\arraybackslash}p{0.19\textwidth}
>{\raggedright\arraybackslash}p{0.19\textwidth}
>{\raggedright\arraybackslash}p{0.25\textwidth}
>{\raggedright\arraybackslash}p{0.25\textwidth}@{}}
\caption{The first-build risk register.}
\label{tab:first-build-risks}\\
\toprule
Risk and owner & Early signal & Mitigation & Stop or redirect trigger \\
\midrule
Silent semantic change; product engineer and domain reviewer & A seeded sign,
unit, citation, or scope change receives a clean comparison. & Typed fixtures,
dual representation, and explicit abstention. & Any unexplained material
change blocks the live path. \\
False-negative preflight; evaluation lead & A seeded dependency, evidence, or
scope defect reaches a release packet, or a critic agrees with a bad draft. &
Held-out mutation suite, isolated critics, frozen thresholds, and fail-closed
unknown state. & Any critical mutation can pass without a recorded abstention. \\
Repair oscillation; product engineer & A repair fixes one gate and reopens the
same gate or alternates between two failures. & Cause-specific strategies,
parent hashes, cycle limit, and oscillation detector. & The controller loops or
silently hands the loop to the reader. \\
False-alert burden; evaluation lead & Held-out triage time meets or exceeds
incumbent issue-finding time. & Remove or narrow the rule family; preserve the
author's rejection reason. & No diagnostic family earns positive net attention
on held-out review. \\
Friendly-case selection; evaluation lead & Eligibility exclusions concentrate
on difficult genres, authors, or privacy classes. & Freeze eligibility and keep
all exclusions with reasons. & The selected case cannot represent the stated
beachhead; report feasibility only and redesign sampling. \\
Prior exposure and reader contamination; independent coding reviewer & A
reader recognizes the project history or earlier version. & Record exposure,
use a fresh reader where possible, and report a flagged sensitivity analysis. &
No uncontaminated reader is available for the declared task. \\
Coaching or unblinding; evaluation lead & One arm receives extra explanation or
the interface reveals its workflow. & Script equal context, freeze answer keys,
and preserve session logs. & Assignment cannot be concealed or equivalent
context cannot be supplied. \\
Nonacceptance, decline, or abandonment; evaluation lead & Documents disappear
from the dataset after a short or failed path. & Fixed horizon and explicit
terminal dispositions for every assigned case. & Outcomes cannot be recovered
for a material share of assigned documents. \\
Privacy breach; technical owner & An undeclared network call or full passage
appears in a log. & Network-deny test, field-level approval, minimization, and
incident response. & Any unauthorized external transmission stops the run. \\
Parser or dependency failure; product engineer & Unknown macros receive a
false pass or an update changes locations. & Pin versions, keep a fallback,
replay fixtures, and return to the incumbent path. & Neither adapter can locate
or explicitly abstain on the protected object set. \\
Scope overrun; sponsor & Optional model, Markdown, drift, or editor work delays
the LaTeX vertical slice. & Weekly critical-path review and deletion of deferred
features. & Brief, plan, draft, preflight, and release are not usable at the end
of week 5. \\
Incomplete scholarship; research lead & Load-bearing rows remain singly
screened or abstract-only. & Independent screening, full-text anchors,
design-specific appraisal, and external review. & Strong literature or
efficacy language remains blocked; implementation learning may continue. \\
\bottomrule
\end{longtable}

\section{The choice after twelve weeks}

The recommendation is to fund the bounded build, corpus, and one feasibility
document. At week twelve the sponsor should receive enough evidence to choose
among three paths:

\begin{table}[H]
\centering
\small
\begin{tabularx}{\textwidth}{@{}p{0.2\textwidth}L p{0.28\textwidth}@{}}
\toprule
Choice & Evidence required & What happens next \\
\midrule
Proceed & The packet is usable, protected objects are preserved, reader burden is
credible, and the comparative endpoint is observable. & Preregister and fund
the paired or stepped comparative study. \\
Revise & The product is technically viable but one diagnostic, parser, or reader
interaction adds avoidable burden. & Repair the named component and repeat the
bounded feasibility test. \\
Stop & Meaning changes are unexplained, privacy is unacceptable, or the reader
still cannot recover purpose and action. & Preserve the evidence and return the
remaining person-weeks and budget. \\
\bottomrule
\end{tabularx}
\caption{The build earns a larger investment only by changing this decision
from a guess into an observed choice.}
\label{tab:authorization-outcomes}
\end{table}

\begin{decisionbox}{What this investment covers}
Approve the twelve-week MVP and feasibility run described here. Deployment,
per-document authorship detection, autonomous acceptance, and any claim that
Humanvoice improves expert writing remain decisions for the sponsor after the
team has observed the result and priced the comparison.
\end{decisionbox}

The first investment is intentionally narrow, but it is not the end of the
argument. A feasibility run can establish that the path is executable and that
its records are interpretable; it cannot establish that the same path is
valuable across writers, genres, or readers. The next chapter treats that
distinction as a research design problem. It specifies the measurement ladder,
the held-out corpus, and the comparative program that must follow a successful
build. In this order, implementation earns the right to ask a larger question
without quietly answering it in advance.
